\documentclass[10pt]{article}
\usepackage[T1]{fontenc}
\usepackage[utf8]{inputenc}
\usepackage[french]{babel}
\usepackage{amsmath,amssymb,mathrsfs,mathpazo}
\usepackage{geometry}
\geometry{a4paper,margin=2cm}
\usepackage{graphicx,array,tabularx,tabularray,colortbl,multirow,multicol,booktabs,enumitem,pifont,ulem,fancybox,diagbox,cancel}
\usepackage{pstricks,pstricks-add,pst-plot,pst-tree,pst-node,pst-text,pst-func,pst-math,pst-bspline,pst-3dplot,pst-eucl,pst-solides3d}
\let\pscalculate\undefined
\usepackage{pst-bezier}
\usepackage{tikz,pgfplots,tkz-tab}
\usetikzlibrary{arrows,arrows.meta,patterns,calc,positioning,trees,decorations.markings,shapes.geometric,shapes.misc,shapes.arrows,intersections,decorations.pathmorphing,decorations.pathreplacing,backgrounds,shadows,babel}
\pgfplotsset{compat=1.18}
\usepgfplotslibrary{fillbetween}
\usepackage[np]{numprint}
\usepackage{listings,makecell,pas-tableur,twoopt}
\let\barre\undefined
\usepackage[pstricks]{tablvar}
\definecolor{darkgreen}{RGB}{0,100,0}
\providecommand{\red}{\color{red}}
\providecommand{\blue}{\color{blue}}
\providecommand{\green}{\color{green}}
\providecommand{\black}{\color{black}}
\providecommand{\white}{\color{white}}
\providecommand{\gray}{\color{gray}}
\providecommand{\cyan}{\color{cyan}}
\providecommand{\magenta}{\color{magenta}}
\providecommand{\yellow}{\color{yellow}}
\providecommand{\violet}{\color{violet}}
\providecommand{\orange}{\color{orange}}
\providecommand{\euro}{\texteuro}
\providecommand{\EUR}{\texteuro}
\providecommand{\ieme}{\textsuperscript{e}}
\providecommand{\degres}{\ensuremath{^\circ}}
\providecommand{\degree}{\ensuremath{^\circ}}
\providecommand{\vv}[1]{\overrightarrow{#1}}
\providecommand{\Coord}[1]{#1}
\providecommand{\decosix}{}
\providecommand{\starredbullet}{\ensuremath{\bullet}}
\pagestyle{empty}
\setlength{\parindent}{0pt}
\usepackage{adjustbox}\DeclareUnicodeCharacter{2236}{:}\DeclareUnicodeCharacter{2190}{\ensuremath{\leftarrow}}
\providecommand{\R}{\mathbb{R}}
\providecommand{\N}{\mathbb{N}}
\providecommand{\Z}{\mathbb{Z}}
\providecommand{\Q}{\mathbb{Q}}
\providecommand{\D}{\mathbb{D}}
\providecommand{\C}{\mathbb{C}}
\providecommand{\e}{{\mathrm{e}}}
\providecommand{\vect}[1]{\overrightarrow{#1}}
\providecommand{\vectt}[1]{\overrightarrow{\mathrm{#1}}}
\providecommand{\vvt}[1]{\overrightarrow{\mathrm{#1}}}
\providecommand{\np}[2][]{{#2~#1}}
\providecommand{\vv}[1]{{\overrightarrow{#1}}}
\renewcommand{\d}{{\mathrm{d}}}
\providecommand{\degres}{{^\circ}}
\providecommand{\degre}{{^\circ}}
\providecommand{\blue}{\color{blue}}
\providecommand{\red}{\color{red}}
\providecommand{\enskip}{\quad}
\providecommand{\llbracket}{\lbrack\!\lbrack}
\providecommand{\rrbracket}{\rbrack\!\rbrack}
\providecommand{\euro}{{\text{€}}}
\providecommand{\no}{{\text{n°}}}
\providecommand{\interval}[2]{{[#1;#2]}}
\providecommand{\intervFO}[2]{{[#1;#2[}}
\providecommand{\dotfill}{\ldots}

\providecommand{\ds}{\displaystyle}
\providecommand{\textding}[1]{#1}
\providecommand{\Oij}{\(\mathrm{O};\vec{\imath},\vec{\jmath}\)}
\providecommand{\Oijk}{\(\mathrm{O};\vec{\imath},\vec{\jmath},\vec{k}\)}
\providecommand{\Ouv}{\(\mathrm{O};\vec{u},\vec{v}\)}
\providecommand{\intervFF}[2]{\ensuremath{[#1;#2]}}
\renewcommand{\euro}{\ifmmode\text{\texteuro}\else\texteuro\fi}
\DeclareUnicodeCharacter{00B0}{\ifmmode^{\circ}\else\textdegree\fi}
\DeclareUnicodeCharacter{20AC}{\ifmmode\text{\texteuro}\else\texteuro\fi}
% Énoncés : Éducation nationale. Transcription et corrigés : APMEP, Denis Vergès et contributeurs. Découpage et classement : CoopMaths, Sébastien Lozano, Éric Elter et contributeurs. Adaptation de la présentation : Maths Explorer.
% CC BY-SA : https://coopmaths.fr/www/about/
% https://coopmaths.fr/alea/static/bac/2025/tex/bac_2025_06_sujet2_metropole_3.tex
\begin{document}
\noindent\textit{APMEP / CoopMaths. CC BY-SA. Présentation adaptée par Maths Explorer.}\par\medskip
\medskip

\emph{Pour chacune des affirmations suivantes, indiquer si elle est vraie ou fausse.\\
Justifier chaque réponse. Une réponse non justifiée ne rapporte aucun point.}

\medskip

\begin{enumerate}
\item La suite $(u_n)$ est définie pour tout entier naturel $n$ par
\[u_n = \dfrac{1 + 5^n}{2 +3^n}.\]

\textbf{Affirmation 1 :} La suite $(u_n)$ converge vers $\dfrac 53$.
\item On considère la suite $(w_n)$ définie par:

\begin{center}$w_0 = 0$\: et, pour tout entier naturel\: $n, w_{n+1} = 3w_n - 2n + 3$.\end{center}

\textbf{Affirmation 2 :} Pour tout entier naturel $n,\: w_n \geqslant n$.
\item On considère la fonction $f$ définie sur $]0~;~+\infty[$ dont la courbe représentative $\mathcal{C}_f$ est donnée dans un repère orthonormé sur la figure (Fig. 1) en page suivante.

On précise que:

\begin{itemize}[label=$\bullet~$]
\item $T$ est la tangente à $\mathcal{C}_f$ au point A d'abscisse 8 ;
\item L'axe des abscisses est la tangente horizontale à $\mathcal{C}_f$ au point d'abscisse 1.
\end{itemize}
\begin{center}
\psset{unit=1cm,arrowsize=2pt 3}
\begin{pspicture*}(-0.5,-1.5)(12,8)
\psaxes[linewidth=1.25pt,Dx=15,Dy=15]{->}(0,0)(0,0)(12,8)
\psgrid[gridlabels=0pt,subgriddiv=1,gridwidth=0.2pt](0,0)(12,8)
\psplot[plotpoints=2000,linewidth=1.25pt,linecolor=blue]{0.01}{12}{x 0.0195 exp 1 sub 50 mul x ln  mul}
\uput[dr](11,5.6){\blue $\mathcal{C}_f$} \uput[ul](11,6){$T$}
\uput[u](8,4.3026){A}\psdot(8,4.3026)
\psline(-0.3,0)(8,4.3026)(16.3,8.6052)
\rput(6,-0.5){\emph{Fig. 1}}
\uput[d](1,0){1}\uput[l](0,1){1}\uput[dl](0,0) {O}
\end{pspicture*}
\end{center}

\textbf{Affirmation 3 :} D'après le graphique, la fonction $f$ est convexe sur son ensemble
de définition.

\item  \textbf{Affirmation 4 :} Pour tout réel $x > 0,\quad \ln (x) - x + 1 \leqslant 0$, où $\ln$ désigne la fonction logarithme népérien.
\end{enumerate}

\bigskip

\end{document}
